% Pursuing Stacks (cotes 134-2 to 134-8): archive cote 134-6, site page 65 (PDF page 66), typescript page 383.
% typescript. Draft compared: ps5.tex:4658-4716.
% AI-transcribed by Claude Opus 5.5 (claude-opus-5-5, adaptive thinking, effort high), Message Batches API, batch ps3 (msgbatch_01XdFmuVVBAcgc7EFxdtCqsj), 2026-11-12; not yet checked by eye.
\scanpage{134-6/65}
\mspage{383}
more familiar (to me) cohomology pot ! Therefore I'm not going to write
out this proof, as ``the'' natural proof is going to come out by itself,
once we got a good conceptual understanding \xout{what}\inserted[typed]{of} homology,\xout{and} cohomology
and abelianization, in the context of \xout{the modelizing (Cat)} ``spaces''
embodied by objects of (Cat). Thus, I feel what is mainly needed now\hand{/}is
an overall review of the relevant notions and facts along these lines -
most of which we've come in touch with before, be it only ``en passant''.

Before starting, just an afterthought on terminology. It occured
to me that the name of an ``integrator''\inserted[typed]{(for A)}\hand{,} for a projective resolution of
the constant abelian presheaf $\underline{Z}_A$ in $A^{\wedge}_{ab}$\hand{,} is unaccurate - as it was meant
to suggest that it's main use is for allowing computation,\xout{of} for an
arbitrary abelian presheaf (or complex of such presheaves) F on A, of
$R\Gamma(A,F)$, which we were thinking of by that time as the ``integration''
of F over A (or over the associated topos). But it has turned out that
for the sake of coherence with a broader use of the notions of ``integra-%
tion'' and ``cointegration'' (compare section 69), the appropriate designa-%
tion of $R\Gamma(A,F)$ is ``\uline{co}integration'' of F over A, not integration. There-%
fore the appropriate designation for a projective resolution $L^A_{\bullet}$ of $\underline{Z}_A$ ,
allowing computational\xout{\unsure{l}} expression of cointegration, is ``\uline{cointegrator}''
\inserted[typed]{(for A)} rather than``integrator''. On the other hand, in terms of the dual catego-%
ry $B=A^{o}$ , it turns out that such $L^A_{\bullet}$ allows computational expression of
\uline{integration} (i.e. homology) over B, and therefore it seems adequate to
call $L^A_{\bullet}$ also an ``\uline{integrator}'' \uline{for} B . Moreover, it turns out that such
an integrator for B is equally an ``abelianizator'' for B, i.e. it allows
simultaneous computation of the ``abelianizations'' of the homotopy types
defined\xout{s} by \xout{an} arbitrary object\hand{s} X in $A^{\wedge}$ , in terms of the abelianization
$Wh_A(X) = \underline{Z}^{(X)}$ of X (cf sections 93 and 102) - and possibly the converse
holds too. Whether this is so or not, there does'nt seem at present
much sense to bother about abelianizators which do not come from integra-%
tors, while the latter have the unvaluable advantage \xout{of being} (besides
mere existence) of being unique up to homotopism. Thus in practical
terms, it would seem that abelianizators (for a given small category B)
are no more no less than just integrators (for the same B, i.e. cointe-%
grators for $A = B^{o}$) - and I would therefore suggest to simply drop the
designation ``abelianizator'' for the benefit of the synonym\xout{\unsure{e}} ``integrator'',
which fits more suggestively into the pair of dual notions integrator -
cointegrator.

\margin{65}
